% !TeX root = 系统动力学建模与仿真笔记.tex
本笔记是《装备系统动力学建模与仿真》课程的学习总结，内容主要依据 Thomas R. Kane 的《Spacecraft Dynamics》一书和张建书老师的讲义整理而成，涉及刚体的姿态描述（方向余弦、欧拉参数、罗德里格斯参数、角速度矩阵与角速度矢量等）及其在运动学分析中的应用。为使全文记号统一，先给出本书采用的符号约定说明，后续各章均遵循本节约定；详细的编辑与规范化要求另见随笔记维护的《符号约定与工作要求》文档。

\section{符号说明}
\label{sec:符号说明}

本书采用如下字型约定：\textbf{向量与矩阵一律用粗体斜体（$\bm{}$）表示}，例如 $\bm{a}$、$\bm{A}$、$\bm{\omega}$；\textbf{标量、分量与矩阵元素不加粗}，例如 $a_i$、$a_{ij}$、$\omega_i$。单位向量组用 $\mathbf{}$（正粗体）表示，如 $\mathbf{a}_i$、$\mathbf{b}_i$。凡表示"在参考系 $x$ 中表达或取分量"的下标一律写成 $|x$，如 $\nu_{i|A}$、$\bm{\nu}_{|A}$。


\begin{longtable}{ll}
\caption{数学基础记号}\label{tab:符号说明:数学记号}\\
\toprule
符号 & 含义 \\
\midrule
\endfirsthead
\multicolumn{2}{c}{\tablename~\thetable{}（续表）}\\
\toprule
符号 & 含义 \\
\midrule
\endhead
$\mathbb{R}^{n}$ & $n$ 维实向量空间 \\
$\bm{A}^{\mathrm{T}}$ & 矩阵 $\bm{A}$ 的转置 \\
$\bm{u}\cdot\bm{v}$ & 向量 $\bm{u}$ 与 $\bm{v}$ 的点积 \\
$\bm{u}\times\bm{v}$ & 向量 $\bm{u}$ 与 $\bm{v}$ 的叉积 \\
$\tilde{\bm{u}}$ & 向量 $\bm{u}$ 的叉乘矩阵，满足 $\tilde{\bm{u}}\bm{v}=\bm{u}\times\bm{v}$ \\
$\delta_{ij}$ & Kronecker 符号 \\
$\epsilon_{ijk}$ & Levi--Civita 符号 \\
$\mathrm{tr}(\bm{A})$ & 矩阵 $\bm{A}$ 的迹 \\
$\det(\bm{A})$ & 矩阵 $\bm{A}$ 的行列式 \\
$\left|\bm{v}\right|$ & 向量 $\bm{v}$ 的模长 \\
$\dfrac{\mathrm{d}}{\mathrm{d}t}$ & 对时间 $t$ 求导 \\
\bottomrule
\end{longtable}

\begin{longtable}{ll}
\caption{参考系与单位向量记号}\label{tab:符号说明:参考系}\\
\toprule
符号 & 含义 \\
\midrule
\endfirsthead
\multicolumn{2}{c}{\tablename~\thetable{}（续表）}\\
\toprule
符号 & 含义 \\
\midrule
\endhead
$A$,\, $B$ & 参考系（或刚体） \\
$\mathbf{a}_1,\mathbf{a}_2,\mathbf{a}_3$ & 固定在 $A$ 中的右手正交单位向量组 \\
$\mathbf{b}_1,\mathbf{b}_2,\mathbf{b}_3$ & 固定在 $B$ 中的右手正交单位向量组 \\
$\mathbf{e}_{i|A}$ & 在参考系 $A$ 中表达的单位向量 $\mathbf{e}_i$ \\
$\mathbf{v}_{|A}$ & 向量 $\bm{v}$ 在参考系 $A$ 中的投影列阵 \\
$\bm{\nu}_{|A}$ & 向量 $\bm{v}$ 在参考系 $A$ 中的分量行阵 \\
$\nu_{i|A}$ & 分量 $\bm{\nu}_{|A}$ 的第 $i$ 个元素 \\
$\dfrac{\mathrm{d}\bm{v}}{\mathrm{d}t}\Big|_{A}$ & 向量 $\bm{v}$ 在参考系 $A$ 中的时间导数 \\
\bottomrule
\end{longtable}

\begin{longtable}{ll}
\caption{转动描述记号}\label{tab:符号说明:转动}\\
\toprule
符号 & 含义 \\
\midrule
\endfirsthead
\multicolumn{2}{c}{\tablename~\thetable{}（续表）}\\
\toprule
符号 & 含义 \\
\midrule
\endhead
$\bm{\lambda}$ & 转轴方向的单位向量 \\
$\theta$ & 转角 \\
$\bm{A}$ & 方向余弦矩阵 \\
$a_{ij}$ & 方向余弦矩阵的元素，$a_{ij}=\mathbf{e}_{i|A}\cdot\mathbf{e}_{j|B}$ \\
$\bm{I}$ & 单位矩阵 \\
$\bm{\epsilon}$ & 欧拉参数向量部分，$\bm{\epsilon}=\bm{\lambda}\sin\frac{\theta}{2}$ \\
$\epsilon_{4}$ & 欧拉参数标量部分，$\epsilon_{4}=\cos\frac{\theta}{2}$ \\
$\bm{\rho}$ & 罗德里格斯向量，$\bm{\rho}=\bm{\lambda}\tan\frac{\theta}{2}$ \\
$\bm{\omega}$ & 角速度向量 \\
$\omega_i$ & 角速度分量 $(i=1,2,3)$ \\
$\bm{\tilde{\omega}}$ & 角速度矩阵，$\bm{\tilde{\omega}}=\bm{A}^{\mathrm{T}}\dot{\bm{A}}$ \\
\bottomrule
\end{longtable}

\begin{longtable}{ll}
\caption{相对转动量记号}\label{tab:符号说明:相对量}\\
\toprule
符号 & 含义 \\
\midrule
\endfirsthead
\multicolumn{2}{c}{\tablename~\thetable{}（续表）}\\
\toprule
符号 & 含义 \\
\midrule
\endhead
$\bm{\omega}_{B|A}$ & $B$ 在 $A$ 中的角速度（竖线前为物体，竖线后为参考系） \\
$\bm{\omega}_{A|B}$ & $A$ 在 $B$ 中的角速度 \\
$\bm{\omega}_{B|A}$ & 反对称性：$\bm{\omega}_{B|A}=-\bm{\omega}_{A|B}$ \\
$\bm{\rho}_{AB}$ & $B$ 相对 $A$ 的罗德里格斯向量 \\
$\rho_{AB,i}$ & 罗德里格斯向量 $\bm{\rho}_{AB}$ 的第 $i$ 个分量 \\
$\bm{\sigma}$ & 单位并矢 \\
\bottomrule
\end{longtable}

\begin{longtable}{ll}
\caption{运动学与动力学量记号}\label{tab:符号说明:运动学}\\
\toprule
符号 & 含义 \\
\midrule
\endfirsthead
\multicolumn{2}{c}{\tablename~\thetable{}（续表）}\\
\toprule
符号 & 含义 \\
\midrule
\endhead
$\mathbf{p}$ & 位置向量 \\
$\bm{\delta}$ & 位移矢量 \\
$\bm{\delta}^{*}$ & 螺杆平移矢量 \\
$s$ & 弧长位移 \\
$\rho$ & 主曲率半径 \\
$\lambda$ & 空间曲线的挠率 \\
$\bm{I}_{ij|A}$ & 惯性并矢 $\bm{I}$ 在参考系 $A$ 中的分量 \\
\bottomrule
\end{longtable}

\begin{mydefinition}{爱因斯坦求和约定}
    在同一项中，若某一角标出现两次，则隐含对该角标在约定范围内求和，本书中求和范围默认为 $i,j,k,\dots=1,2,3$。出现两次的角标称为\emph{哑指标}（dummy index），求和后该指标消失；只出现一次的角标称为\emph{自由指标}（free index），它代表其取值范围内的任一值，等式两端自由指标必须一致。
\end{mydefinition}

举一个最简单的例子。设 $\bm{a}$、$\bm{b}$ 为三维向量，按求和约定，
\begin{equation}
    a_ib_i=a_1b_1+a_2b_2+a_3b_3=\bm{a}\cdot\bm{b},
    \label{eq:符号说明:求和内积}
\end{equation}
即 $a_ib_i$ 正是两向量的内积。作为具体的数值例子，若 $\bm{a}=(1,2,3)$、$\bm{b}=(4,5,6)$，则
\begin{equation}
    a_ib_i=1\cdot4+2\cdot5+3\cdot6=4+10+18=32.
\end{equation}

矩阵乘向量同样可以简洁地写出：设 $\bm{A}$ 的元素为 $a_{ij}$，则
\begin{equation}
    (\bm{A}\bm{x})_i=a_{ij}x_j=a_{i1}x_1+a_{i2}x_2+a_{i3}x_3.
\end{equation}

叉积借助 Levi--Civita 符号可以写成
\begin{equation}
    (\bm{u}\times\bm{v})_i=\epsilon_{ijk}u_jv_k,
\end{equation}
例如取 $i=1$ 时，由于 $\epsilon_{ijk}$ 仅在下标为循环排列 $(1,2,3)$ 时取 $1$、逆循环排列时取 $-1$，其余为零，故只有 $\epsilon_{123}=1$、$\epsilon_{132}=-1$ 两项非零，
\begin{equation}
    (\bm{u}\times\bm{v})_1=\epsilon_{123}u_2v_3+\epsilon_{132}u_3v_2=u_2v_3-u_3v_2,
\end{equation}
与叉积的定义一致。

在本笔记中这一约定大量出现，例如角速度矩阵一节的泊松运动学方程
\begin{equation}
    \dot{a}_{ij}=a_{ig}\,\omega_h\,\epsilon_{ghj}
\end{equation}
中，$g,h$ 为哑指标（求和），$i,j$ 为自由指标。此外，哑指标的具体名称无关紧要，例如 $a_ib_i=a_jb_j$，这一点在多重求和时尤为方便。
